% =====================================================================
% Review: Adsorption Kinetic Models (liquid-phase batch systems)
% Narrative review (English); target: international journal (TBD)
% Angle: A (practitioner's workflow) + B (unifying rate-limiting-step
%        framework) + reproducible re-analysis of published datasets
% Build: latexmk -pdf main.tex   (see ../README.md)
% =====================================================================
\documentclass[11pt,a4paper]{article}

\usepackage[T1]{fontenc}
\usepackage[utf8]{inputenc}
\usepackage{lmodern}
\usepackage[american]{babel}
\usepackage[margin=2.5cm]{geometry}
\usepackage{amsmath,amssymb}
\usepackage{siunitx}
\usepackage{booktabs}
\usepackage{graphicx}
\usepackage{csquotes}
\usepackage[style=apa,backend=biber]{biblatex}
\addbibresource{references.bib}
\usepackage[hidelinks]{hyperref}

% Placeholder for unwritten content
\newcommand{\todo}[1]{\textbf{[TODO: #1]}}

\title{Adsorption Kinetics in Liquid-Phase Batch Systems: From Rate-Limiting
       Steps to Robust Model Selection}
\author{Author Name\thanks{Affiliation; corresponding e-mail.}}
\date{}

\begin{document}
\maketitle

% ---------------------------------------------------------------------
\begin{abstract}
\todo{Abstract, 200--250 words: (1) kinetic models are routinely fitted
but rarely interpreted correctly; (2) this review organizes liquid-phase
batch models by rate-limiting step and shows how they are related;
(3) a workflow for design, fitting, and selection, demonstrated by
re-analysis of published datasets; (4) reporting checklist.}

\medskip
\noindent\textbf{Keywords:} adsorption kinetics; batch adsorption;
pseudo-second order; intraparticle diffusion; non-linear regression;
model selection
\end{abstract}

% =====================================================================
\section{Introduction}
\label{sec:intro}

\todo{P1: Adsorption is a leading option for removing dyes, metals, and
emerging contaminants from water; low-cost adsorbents dominate the
literature \parencite{Bailey1999_LowCostHeavyMetals,Crini2006_LowCostDye}.}

\todo{P2: Isotherms describe equilibrium \parencite{Foo2010_IsothermModeling};
kinetics describe rate and, potentially, mechanism. Kinetic parameters feed
contact-time selection and reactor design.}

\todo{P3: The problem --- PSO is reported as ``best fit'' in most studies,
often from linearized plots, and then used to claim chemisorption
\parencite{Tran2017_AdsorptionMistakes,Simonin2016_PFOvsPSO}.}

\todo{P4: Existing reviews catalog models \parencite{Wang2020_KineticModels}
or list mistakes \parencite{Tran2017_AdsorptionMistakes}; none links the
physical basis of each model to a practical, tested workflow.}

\todo{P5: Aim, scope (liquid-phase, batch, single solute; column and
gas-phase systems excluded), and organization.}

% =====================================================================
\section{Physical Basis of Batch Adsorption Kinetics}
\label{sec:basis}

\subsection{Mass balance in a batch reactor}
The amount adsorbed at time $t$ follows from a mass balance:
\begin{equation}
  q_t = \frac{(C_0 - C_t)\,V}{m},
  \label{eq:qt}
\end{equation}
where $q_t$ is the amount adsorbed (\si{\milli\gram\per\gram}), $C_0$ and
$C_t$ are the initial and time-$t$ concentrations
(\si{\milli\gram\per\liter}), $V$ is the solution volume (\si{\liter}),
and $m$ is the adsorbent mass (\si{\gram}).

\todo{Emphasize that $C_t$ decreases during the run (finite bath), which
many models implicitly ignore.}

\subsection{Resistances in series}
\todo{Figure 1: schematic of a particle with (i) bulk transport,
(ii) external film diffusion, (iii) intraparticle (pore/surface)
diffusion, (iv) attachment at the active site. The slowest step controls
the overall rate.}

\subsection{Diagnosing the controlling step}
\todo{Effect of agitation speed (film), particle size (intraparticle),
and $C_0$ and temperature on apparent rate constants; Biot number for
relative film vs intraparticle resistance.}

% =====================================================================
\section{Surface-Reaction Models}
\label{sec:reaction}

\subsection{Langmuir kinetics as the parent model}
Assuming reversible attachment on a homogeneous monolayer,
\begin{equation}
  \frac{dq_t}{dt} = k_a C_t (q_m - q_t) - k_d q_t,
  \label{eq:langmuir-kin}
\end{equation}
where $k_a$ and $k_d$ are the adsorption and desorption rate constants and
$q_m$ is the monolayer capacity. At equilibrium Eq.~\eqref{eq:langmuir-kin}
reduces to the Langmuir isotherm. If $C_t \approx C_0$ (large excess of
adsorbate), Eq.~\eqref{eq:langmuir-kin} becomes linear in $q_t$ and yields
the PFO form (Section~\ref{sec:pfo}) with $k_1 = k_a C_0 + k_d$.

\todo{Cite the analysis showing PFO at high $C_0$ and PSO-like behavior at
low $C_0$ (Azizian, 2004 --- not yet verified).}

\subsection{Pseudo-first order (PFO)}
\label{sec:pfo}
\begin{equation}
  \frac{dq_t}{dt} = k_1 (q_e - q_t)
  \quad\Longrightarrow\quad
  q_t = q_e\left(1 - e^{-k_1 t}\right),
  \label{eq:pfo}
\end{equation}
with the linear form $\ln(q_e - q_t) = \ln q_e - k_1 t$, which requires
$q_e$ to be known in advance.
\todo{Cite original Lagergren (1898) --- not yet verified.}

\subsection{Pseudo-second order (PSO)}
\label{sec:pso}
\begin{equation}
  \frac{dq_t}{dt} = k_2 (q_e - q_t)^2
  \quad\Longrightarrow\quad
  q_t = \frac{k_2 q_e^2\, t}{1 + k_2 q_e\, t},
  \label{eq:pso}
\end{equation}
whose most widely used linear form is
\begin{equation}
  \frac{t}{q_t} = \frac{1}{k_2 q_e^2} + \frac{t}{q_e},
  \label{eq:pso-lin}
\end{equation}
and the initial adsorption rate is $h = k_2 q_e^2$
\parencite{Ho1999_PseudoSecondOrder}.

\todo{Note that $t/q_t$ vs $t$ is dominated by late-time points, so a high
$R^2$ is almost guaranteed \parencite{Simonin2016_PFOvsPSO}.}

\subsection{Mixed-order and \texorpdfstring{$n$th}{nth}-order models}
\todo{General $n$th-order and mixed first/second-order rate laws that
contain PFO and PSO as limits (sources not yet verified).}

\subsection{Empirical models: Elovich and Avrami}
\begin{equation}
  \frac{dq_t}{dt} = \alpha\, e^{-\beta q_t}
  \quad\Longrightarrow\quad
  q_t = \frac{1}{\beta}\ln(1 + \alpha\beta t),
  \label{eq:elovich}
\end{equation}
\begin{equation}
  q_t = q_e\left\{1 - \exp\!\left[-(k_{AV}\, t)^{n_{AV}}\right]\right\}.
  \label{eq:avrami}
\end{equation}
\todo{Interpretation: Elovich for energetically heterogeneous surfaces;
Avrami fractional order $n_{AV}$. Stress that both are descriptive, not
mechanistic. Original sources not yet verified.}

% =====================================================================
\section{Diffusion Models}
\label{sec:diffusion}

\subsection{External film diffusion}
Transfer across the liquid film is commonly written as
\begin{equation}
  \frac{dC_t}{dt} = -\beta_L S\,(C_t - C_s),
  \label{eq:film}
\end{equation}
where $\beta_L$ is the external mass-transfer coefficient
(\si{\centi\meter\per\second}), $S$ the external particle surface area
per unit solution volume (\si{\per\centi\meter}), and $C_s$ the
concentration at the particle surface. At early times $C_s \approx 0$, so
$\ln(C_t/C_0) \approx -\beta_L S\, t$.
\todo{Citation for the initial-rate method --- not yet verified.}

\subsection{Boyd film model}
With $F = q_t/q_e$,
\begin{equation}
  F = 1 - \frac{6}{\pi^2}\sum_{n=1}^{\infty}\frac{1}{n^2}
      \exp\!\left(-n^2 B t\right),
  \label{eq:boyd}
\end{equation}
which is approximated by
\begin{equation}
  Bt =
  \begin{cases}
    \left(\sqrt{\pi} - \sqrt{\pi - \dfrac{\pi^2 F}{3}}\right)^{2},
      & F \le 0.85,\\[2ex]
    -0.4977 - \ln(1 - F), & F > 0.85.
  \end{cases}
  \label{eq:boyd-approx}
\end{equation}
A $Bt$ versus $t$ line that does not pass through the origin indicates
that film diffusion contributes to rate control.
\todo{Cite Boyd et al.\ (1947) --- not yet verified.}

\subsection{Intraparticle diffusion (Weber--Morris)}
\begin{equation}
  q_t = k_{id}\, t^{1/2} + C,
  \label{eq:weber-morris}
\end{equation}
where $C$ relates to the boundary-layer thickness. A multilinear plot
indicates more than one rate-controlling step.
\todo{Cite Weber \& Morris (1963) --- not yet verified. Discuss the
subjectivity of choosing breakpoints.}

\subsection{Homogeneous surface diffusion model (HSDM)}
For a spherical particle of radius $R$,
\begin{equation}
  \frac{\partial q}{\partial t} =
  \frac{D_s}{r^2}\frac{\partial}{\partial r}
  \left(r^2 \frac{\partial q}{\partial r}\right),
  \qquad 0 < r < R,
  \label{eq:hsdm}
\end{equation}
coupled to the film equation at $r = R$ and to the batch mass balance.
\todo{Brief: physically meaningful $D_s$, but requires numerical solution;
point to the code accompanying Section~\ref{sec:reanalysis}.}

% =====================================================================
\section{Linking the Models}
\label{sec:linking}

\todo{Core ``B'' section. (1) PFO as the $C_t \approx C_0$ limit of
Langmuir kinetics (Eq.~\ref{eq:langmuir-kin}); (2) conditions under which
PSO emerges from Langmuir kinetics and from diffusion-controlled uptake;
(3) consequence: goodness of fit alone cannot identify mechanism.}

\todo{Figure 2: map of models as limits of one another.
Figure 3: one synthetic diffusion-controlled curve fitted equally well by
PFO, PSO, and Elovich.}

% =====================================================================
\section{Parameter Estimation and Model Selection}
\label{sec:estimation}

\subsection{Linear vs non-linear regression}
\todo{Linearization changes the error structure and weights points
unequally \parencite{Foo2010_IsothermModeling,Tran2017_AdsorptionMistakes}.
Recommend non-linear least squares on $q_t$ (or $C_t$).}

\subsection{Error functions and information criteria}
\begin{align}
  \mathrm{RMSE} &= \sqrt{\frac{1}{N}\sum_{i=1}^{N}
            (q_{\mathrm{exp}} - q_{\mathrm{cal}})_i^2},
  \label{eq:rmse}\\
  \chi^2 &= \sum_{i=1}^{N}
            \frac{(q_{\mathrm{exp}} - q_{\mathrm{cal}})_i^2}
                 {q_{\mathrm{cal},i}},
  \label{eq:chi2}\\
  \mathrm{AIC_c} &= N \ln\!\left(\frac{\mathrm{SSE}}{N}\right) + 2p
            + \frac{2p(p+1)}{N - p - 1},
  \label{eq:aicc}\\
  \mathrm{BIC} &= N \ln\!\left(\frac{\mathrm{SSE}}{N}\right) + p\ln N,
  \label{eq:bic}
\end{align}
where $N$ is the number of data points, SSE the sum of squared errors,
and $p$ the number of model parameters.
\todo{Why AIC\textsubscript{c} matters: kinetic runs often have
$N \approx 10$--15 points.}

\subsection{Residuals and parameter uncertainty}
\todo{Residual plots, confidence intervals from the covariance matrix or
bootstrap, correlation between $q_e$ and $k$.}

\subsection{Experimental design for kinetics}
\todo{Dense sampling at early times, runs at several $C_0$, agitation and
particle-size tests, replicates.}

\subsection{Recommended workflow}
\todo{Figure 4: decision flowchart from experimental design to reporting.}

% =====================================================================
\section{Re-analysis of Published Datasets}
\label{sec:reanalysis}

\todo{Select 3--5 published liquid-phase batch datasets (different
adsorbate classes: dye, metal, organic micropollutant). Refit all models
with linear and non-linear methods; compare rankings by $R^2$ vs
AIC\textsubscript{c}; report parameter uncertainty. Code in
\texttt{analysis/}.}

\todo{Table 2: dataset sources. Table 3: fitted parameters and criteria.
Figure 5: linear vs non-linear fits for one dataset.}

% =====================================================================
\section{Common Pitfalls and Reporting Checklist}
\label{sec:pitfalls}

\todo{Build on \textcite{Tran2017_AdsorptionMistakes} and
\textcite{Wang2020_KineticModels}: wrong units of $k_2$, calculated vs
experimental $q_e$, too few early-time points, inferring chemisorption from
PSO, ignoring the finite-bath effect.}

\todo{Table 4: reporting checklist for kinetic studies.}

% =====================================================================
\section{Outlook}
\label{sec:outlook}

\todo{Physically based models becoming accessible via open-source solvers;
multicomponent systems; data-driven model selection; standard datasets and
code sharing. Note that column and gas-phase kinetics are outside the
scope of this review.}

% =====================================================================
\section{Conclusions}
\label{sec:conclusions}

\todo{Three to four conclusions and practical recommendations.}

% ---------------------------------------------------------------------
\begin{table}[ht]
  \centering
  \caption{Liquid-phase batch kinetic models covered in this review.}
  \label{tab:summary}
  \begin{tabular}{@{}llll@{}}
    \toprule
    Model & Equation & Parameters & Controlling step \\
    \midrule
    Langmuir kinetics & \eqref{eq:langmuir-kin} & $k_a$, $k_d$, $q_m$
      & Surface reaction \\
    PFO & \eqref{eq:pfo} & $q_e$, $k_1$ & Surface reaction (apparent) \\
    PSO & \eqref{eq:pso} & $q_e$, $k_2$ & Surface reaction (apparent) \\
    Elovich & \eqref{eq:elovich} & $\alpha$, $\beta$ & Empirical \\
    Avrami & \eqref{eq:avrami} & $q_e$, $k_{AV}$, $n_{AV}$ & Empirical \\
    Film diffusion & \eqref{eq:film} & $\beta_L$ & External film \\
    Boyd & \eqref{eq:boyd} & $B$ & External film \\
    Weber--Morris & \eqref{eq:weber-morris} & $k_{id}$, $C$
      & Intraparticle \\
    HSDM & \eqref{eq:hsdm} & $D_s$ & Intraparticle \\
    \bottomrule
  \end{tabular}
\end{table}

\todo{Refer to Table~\ref{tab:summary} in Section~\ref{sec:intro} or
Section~\ref{sec:linking}.}

\printbibliography

\end{document}
